\section{Datasets}
\label{sec:datasets}

\draftnote{DRAFT. Numbers below are final/stable (reports/02d, Phase 1
data-layer notes); prose needs polishing and a proper table generated
by script once finalized.}

We use three real mobility datasets, one of which we split into four
independent metropolitan subsets after finding that a single combined
sample does not provide adequate per-user coverage for our evaluation
protocol (Section~\ref{sec:evaluation-protocol}).

\textbf{GeoLife.} The Microsoft Research GeoLife GPS trajectory
dataset, 181 users, 24.9 million GPS records. Windows are built
respecting trajectory boundaries (a trajectory is a genuine, separate
trip), since GeoLife's \texttt{traj\_id} reflects real trip structure.

\textbf{Foursquare-NYC.} 1,083 users. Unlike GeoLife, check-in data has
no native trajectory identifier; an earlier version of this pipeline
synthesized one from same-calendar-day grouping, which left only 71 of
1,083 users with even one length-10 window. We instead build windows
over each user's full chronological check-in sequence within a split,
ignoring the day boundary.

\textbf{Gowalla (four metro subsets).} Gowalla check-ins subset by
metro area: \texttt{gowalla\_nyc} (1,517 users), \texttt{gowalla\_la}
(1,871), \texttt{gowalla\_chicago} (675), and \texttt{gowalla\_sf\_bay}
(2,797). We treat these as four separate datasets rather than one
pooled Gowalla sample: Gowalla-NYC's per-user window coverage, even
after shortening the window length from 10 to 5 steps, left only 27.9\%
of users with the minimum number of test windows our protocol requires
(Section~\ref{sec:evaluation-protocol}), short of a 50\% target; the
other three metros independently fall in the same 24--37\% range.
Treating all four as separate datasets gives 2,235 usable users across
the four metros combined, where pooling into one relaxed-filter sample
would not have.

\textbf{A note on user-ID sparsity.} Gowalla's \texttt{user\_id} values
are sparse within each metro subset (e.g.\ Gowalla-NYC's maximum id is
7,202 for only 1,517 users with any data, a $4$--$5\times$ inflation).
Naively using the maximum id as the number of classes for a
multiclass classifier -- as our gradient-boosted-tree attacker
architecture does -- made that component of our evaluation
pathologically slow (one fit ran over 14 hours without completing on a
7,202-way problem that is really a 1,517-way one). We remap user ids to
a dense range before any downstream processing; GeoLife and Foursquare
already have dense id spaces and are unaffected.
